\section{Cambio de paradigma en los modelos: de la GPU al Agent}

Esta sección cubre el primer eje: el paradigma de los modelos. La historia empieza con la primera GPU, en 1999. No es que la GPU sea un artículo de investigación, sino que cambió el límite de lo computable. Trabajos como Brook, CUDA, AlexNet, ResNet y Transformer muestran un hecho sencillo: que un paradigma algorítmico se vuelva el eje de una época depende muchas veces de que sepa aprovechar el hardware y los datos.

\begin{figure}[H]
\centering
\includegraphics[width=0.96\textwidth]{model-paradigm-timeline.png}
\caption{Línea de tiempo de los paradigmas de modelos: GPU, CNN, Attention, Transformer, MoE y Agent se relevan uno a otro. Diagrama conceptual propio, elaborado a partir de la conversación entre 00:19:35 y 01:52:58.}
\end{figure}

\begin{knowledgebox}{Lectura de la figura: los paradigmas no surgen de repente}
Antes de AlexNet ya existían la GPU y los datos de imágenes; antes de Transformer existían seq2seq, Attention, las redes residuales y el cómputo en paralelo; antes del Agent existían CoT, la llamada a herramientas y ReAct. Un cambio de paradigma suele ser la conexión repentina de muchos avances acumulados dentro de una ventana determinada de hardware y datos.
\end{knowledgebox}

\subsection{GPU, Brook, AlexNet y la lotería del hardware}

Esta sección examina primero cómo el hardware cambió el destino de los modelos. Brook for GPU representa la exploración temprana del cómputo de propósito general en GPU, mientras que AlexNet, en 2012, combinó la GPU con redes convolucionales profundas para superar ampliamente a todos en ImageNet. El ponente insiste en la «lotería del hardware»: si un método se ajusta precisamente al hardware dominante, se amplifican su velocidad de entrenamiento, su ecosistema de ingeniería y la inversión en investigación; en cambio, una idea puede ser correcta y aun así tener dificultades para volverse el eje de una época.

\begin{figure}[H]
\centering
\includegraphics[width=0.94\textwidth]{hardware-lottery.png}
\caption{Lotería del hardware: las arquitecturas que aprovechan el hardware dominante obtienen un rendimiento compuesto a largo plazo. Diagrama conceptual propio, elaborado a partir de la conversación entre 00:19:35 y 00:30:00, y entre 03:56:38 y 04:03:00.}
\end{figure}

\begin{knowledgebox}{Lectura de la figura: por qué AlexNet marca un parteaguas}
AlexNet no solo «usó una CNN»: también demostró que el aprendizaje profundo, la GPU, el conjunto de datos y las técnicas de ingeniería podían combinarse en una ventaja abrumadora. Derrotó al paradigma dominado por las características diseñadas a mano y anticipó la Bitter Lesson: a medida que el poder de cómputo sigue creciendo, los métodos de aprendizaje generales superan una y otra vez al diseño manual.
\end{knowledgebox}

\subsection{seq2seq, Attention y Transformer}

Lo anterior trató sobre imágenes; esta sección pasa a las secuencias. seq2seq hace que el modelo comprima la secuencia de entrada en un estado y luego la decodifique, mientras que Attention permite que, al decodificar, el modelo atienda directamente a distintas posiciones de la entrada, lo que mitiga los problemas de las oraciones largas y la pérdida de contexto. Transformer va más allá: sustituye el cuello de botella secuencial de las RNN por self-attention y cómputo en paralelo, de modo que puede modelar directamente las relaciones entre los token dentro de la secuencia.

\begin{figure}[H]
\centering
\includegraphics[width=0.96\textwidth]{transformer-origin-ladder.png}
\caption{Del modelado de secuencias a Transformer: seq2seq, Attention, las conexiones residuales y el cómputo en paralelo prepararon el terreno en conjunto. Diagrama conceptual propio, elaborado a partir de la conversación entre 00:34:02 y 01:05:24.}
\end{figure}

\begin{importantbox}{La intuición central de Attention}
Attention permite que el modelo, al procesar la posición actual, calcule directamente su relación con las demás posiciones. Frente a las RNN, que dependen de transmitir el estado oculto paso a paso, Attention acorta el recorrido de la información; frente a la convolución, que necesita varias capas para ampliar el campo receptivo, self-attention puede ver toda la ventana de contexto en una sola capa.
\end{importantbox}

\subsection{ResNet, destilación, AlphaGo Zero y MoE}

Lo anterior trató sobre cómo las arquitecturas procesan imágenes y secuencias; esta sección amplía el alcance a las técnicas de entrenamiento, la transferencia de capacidades y la economía del cómputo. El paradigma de los modelos no es solo la estructura: también incluye el entrenamiento y la transferencia de capacidades. ResNet usa conexiones residuales para resolver la dificultad de entrenar redes más profundas; la destilación plantea que un modelo pequeño puede aprender el conocimiento de un modelo grande; AlphaGo Zero mostró que el aprendizaje por refuerzo y el juego contra sí mismo pueden prescindir de las partidas humanas; y el MoE moderno reduce el costo de activación mediante el enrutamiento entre expertos, de modo que la capacidad de un modelo grande y su costo de cómputo ya no están ligados de forma completamente lineal.

\begin{knowledgebox}{Términos clave: hitos del paradigma de modelos}
\begin{tabularx}{\textwidth}{p{0.20\textwidth}p{0.32\textwidth}X}
\toprule
Hito & Problema que resuelve & Influencia posterior \\
\midrule
ResNet & Degradación y dificultad de entrenamiento de las redes profundas & Lo residual se volvió un componente general de las redes profundas y llegó también a Transformer. \\
Distillation & Cómo transferir las capacidades de un modelo grande a uno pequeño & Sustenta el despliegue en el dispositivo, los modelos estudiante y la compresión de modelos. \\
AlphaGo Zero & Aprender una estrategia sin depender de datos humanos & Las ideas del aprendizaje por refuerzo, el juego contra sí mismo y el test-time search se citan una y otra vez. \\
MoE & Aumentar la capacidad de parámetros controlando el cómputo de cada pasada & Modelos como DeepSeek convirtieron el enrutamiento entre expertos en una vía importante de los LLM modernos. \\
\bottomrule
\end{tabularx}
\end{knowledgebox}

\subsection{CoT, LoRA, ReAct y la transición hacia el Agent}

Esta sección pasa de los modelos base a la forma de usarlos. CoT hizo ver que la entrada del modelo influye en su salida, lo que impulsó la prompt engineering y la context engineering; LoRA generalizó la adaptación de parámetros a bajo costo; ReAct puso el reasoning y el acting en el mismo ciclo, lo que sentó las bases para que el Agent pasara de la teoría a la llamada a herramientas.

\begin{figure}[H]
\centering
\includegraphics[width=0.96\textwidth]{adaptation-agent-stack.png}
\caption{La pila que va de la capacidad a la acción: la destilación, LoRA, CoT y ReAct llevan al modelo hacia un sistema utilizable. Diagrama conceptual propio, elaborado a partir de la conversación entre 01:26:40 y 01:37:10.}
\end{figure}

\begin{knowledgebox}{Lectura de la figura: las capacidades del modelo deben «extraerse» y «recibirse»}
La destilación y LoRA resuelven la transferencia de capacidades y el costo de adaptación; CoT hace que el modelo despliegue su razonamiento de forma explícita, y ReAct conecta ese razonamiento con herramientas externas. El Agent no surgió de la nada: estos mecanismos fueron llevando poco a poco al modelo hacia un ciclo cerrado de acción.
\end{knowledgebox}

\subsection{Bitter Lesson: un baño de realidad para el largo plazo}

Todos los hitos anteriores insinúan una regularidad, y esta sección la expone por completo con The Bitter Lesson. Su tesis es que, a largo plazo, los métodos que dependen del cómputo general y de la búsqueda o el aprendizaje a gran escala suelen superar a los que codifican el conocimiento a mano. El ponente subraya que esto no significa que las características diseñadas a mano sean siempre inútiles, sino que las estructuras manuales suelen funcionar dentro de cierto orden de magnitud de cómputo; una vez que el cómputo y los datos superan un umbral, los métodos de aprendizaje generales vuelven a imponerse.

\begin{figure}[H]
\centering
\includegraphics[width=0.94\textwidth]{bitter-lesson-map.png}
\caption{Bitter Lesson: a largo plazo, el cómputo general y el aprendizaje a escala superan una y otra vez a las estructuras diseñadas a mano. Diagrama conceptual propio, elaborado a partir de la conversación entre 01:45:00 y 01:52:40.}
\end{figure}

\begin{warningbox}{Bitter Lesson no significa «no diseñar estructuras»}
El diseño de estructuras sigue siendo importante, sobre todo cuando hay pocos datos, poco cómputo y un objetivo claro. Pero si una estructura no puede hacer scale junto con el cómputo, los datos y el ecosistema de hardware, podría quedar superada por los métodos generales en la siguiente transición de paradigma.
\end{warningbox}

\subsection{Resumen de la sección}

El eje del paradigma de los modelos nos enseña que el progreso de la IA no es el triunfo de un solo artículo, sino la alineación conjunta del hardware, los datos, los métodos de entrenamiento, la arquitectura y la forma de uso. La GPU hizo viable el aprendizaje profundo, Attention hizo aprendibles las dependencias largas, Transformer aprovechó el hardware paralelo, y CoT/ReAct llevaron al modelo hacia la acción.
